\subsection{媒介変数で表した曲線}
曲線が $x=x(t)$，$y=y(t)$ と表されるとき，$t$ を\textbf{媒介変数}という。両関数が連続な導関数をもち，考えている点で $x'(t)\ne0$ なら，その近くで $x(t)$ は狭義単調であり，$t$ を $x$ の関数として表せる。合成関数と逆関数の微分から
\[
 \frac{\dd y}{\dd x}
 =\frac{y'(t)}{x'(t)}
\]
を得る。$x'(t)=0$ の点ではこの商は使えず，垂直接線などを別に調べる必要がある。さらに二階まで微分できる場合は，$x$ での微分が $\frac1{x'(t)}\frac{\dd}{\dd t}$ に対応するので，
\[
 \frac{\dd^2y}{\dd x^2}
 =\frac1{x'(t)}\frac{\dd}{\dd t}
   \left(\frac{y'(t)}{x'(t)}\right)
\]
となる。$\frac{y''(t)}{x''(t)}$ ではない。

